\subsection{树链剖分 (Heavy-Light Decomposition)}
树链剖分，将树上问题转化为序列问题，常配合线段树使用。
\begin{lstlisting}
ll dfn[100010],de[100010],fa[100010],son[100010];
ll sz[100010],top[100010],rk[100010];
void dfs1(ll u,ll fat){
    fa[u]=fat;de[u]=de[fat]+1;sz[u]=1;son[u]=0;
    for(int i=hd[u];i;i=e[i].nx){
        ll v=e[i].v;if(v==fat)continue;
        dfs1(v,u);
        if(sz[v]>sz[son[u]])son[u]=v;
        sz[u]+=sz[v];
    }
}void dfs2(ll u,ll tp){
    top[u]=tp;dfn[u]=++tim;rk[tim]=u;
    if(son[u])dfs2(son[u],tp);
    for(int i=hd[u];i;i=e[i].nx){
        ll v=e[i].v;if(v==fa[u]||v==son[u])continue;
        dfs2(v,v);
    }
}
ll a[100010];
#define ls(p) (p<<1)
#define rs(p) (p<<1|1)
ll t[100010<<2];
ll tg[100010<<2];
void upd(ll p){t[p]=(t[ls(p)]+t[rs(p)])%mod;}
void build(ll l,ll r,ll p){
    tg[p]=0;
    if(l==r){t[p]=a[rk[l]];return;}
    ll mid=l+r>>1;
    build(l,mid,ls(p));
    build(mid+1,r,rs(p));
    upd(p);
}
void pd(ll p,ll l,ll r){
    if(tg[p]==0)return;ll mid=l+r>>1;
    tg[ls(p)]=(tg[ls(p)]+tg[p])%mod;
    tg[rs(p)]=(tg[rs(p)]+tg[p])%mod;
    t[ls(p)]=(t[ls(p)]+tg[p]*(mid-l+1))%mod;
    t[rs(p)]=(t[rs(p)]+tg[p]*(r-mid))%mod;
    tg[p]=0;
}
void add(ll l,ll r,ll p,ll L,ll R,ll k){
    if(L<=l&&r<=R){t[p]=(t[p]+k*(r-l+1))%mod;tg[p]=(tg[p]+k)%mod;return;}
    pd(p,l,r);
    ll mid=l+r>>1;
    if(mid>=L)add(l,mid,ls(p),L,R,k);
    if(mid<R)add(mid+1,r,rs(p),L,R,k);
    upd(p);
}ll get(ll l,ll r,ll p,ll L,ll R){
    if(L<=l&&r<=R)return t[p];
    pd(p,l,r);
    ll mid=l+r>>1,res=0;
    if(mid>=L)res+=get(l,mid,ls(p),L,R);
    if(mid<R)res+=get(mid+1,r,rs(p),L,R);
    return res%mod;
}
void chad(ll u,ll v,ll k){
    while(top[u]!=top[v]){
        if(de[top[u]]<de[top[v]])swap(u,v);
        add(1,n,1,dfn[top[u]],dfn[u],k);
        u=fa[top[u]];
    }if(dfn[u]>dfn[v])swap(u,v);
    add(1,n,1,dfn[u],dfn[v],k);
}ll chs(ll u,ll v){ll res=0;
    while(top[u]!=top[v]){
        if(de[top[u]]<de[top[v]])swap(u,v);
        res+=get(1,n,1,dfn[top[u]],dfn[u]);
        u=fa[top[u]];
    }if(dfn[u]>dfn[v])swap(u,v);
    res+=get(1,n,1,dfn[u],dfn[v]);
    return res%mod;
}
\end{lstlisting}

\paragraph{多测与注意} 每组读边前清 hd、重置边计数，令 tim=0、sz[0]=de[0]=0，再 dfs1、dfs2、build。son 和 lazy 已在各自流程中重置；若有负增量，按题意归一化取模。
